\documentclass[10pt,a4paper]{amsart}
\usepackage[margin=19mm]{geometry}
\usepackage{amsmath,amssymb,mathtools}
\usepackage[scr=rsfs,cal=euler]{mathalfa}
\usepackage{xfrac}
\usepackage[unicode,hidelinks,bookmarks=false]{hyperref}
\newcommand{\ie}{i.e.\ }
\newcommand{\cf}{cf.\ }
\newcommand{\resp}{respectively\ }
\newcommand{\Subsection}{\subsection}
\providecommand{\nobreakdash}{\nobreak\hbox{-}}
\setlength{\emergencystretch}{2em}
\multlinegap0pt
\usepackage{bm}
\usepackage{bbm}
\usepackage{dsfont}
\usepackage{xspace}
\usepackage{cancel}
\usepackage{mathtools}
\usepackage{cleveref}
\usepackage{amsthm}
\makeatletter
\@ifundefined{theorem}{\newtheorem{theorem}{Theorem}}{}
\@ifundefined{lemma}{\newtheorem{lemma}[theorem]{Lemma}}{}
\makeatother
\newcommand{\Z}{\mathbb Z}
\newcommand{\arr}{\mathrm{arr}}
\newcommand{\dep}{\mathrm{dep}}
\newcommand{\dotcup}{\mathbin{\dot\cup}}
\newcommand{\energyact}{M}               % selected actual energy transfer activities
\newcommand{\wait}{\mathrm{wait}}
\newcommand{\w}{\rho} %energy arc weight for matching purposes
\title[Optimizing Travel Time and Regenerative Energy for Periodic ]{Optimizing Travel Time and Regenerative Energy for Periodic Timetables}
\author{Sarah Roth, Sven J\"ager, Niels Lindner, Anita Sch\"obel}
\date{Source: arXiv:2605.02355v1 (CC BY 4.0)}
\begin{document}
\maketitle

\noindent\emph{Source and licence.} Optimizing Travel Time and Regenerative Energy for Periodic Timetables. Version arXiv:2605.02355v1 (CC BY 4.0), distributed under the Creative Commons Attribution 4.0 International licence, as recorded in the arXiv licence metadata for this exact version and shown on the abstract page https://arxiv.org/abs/2605.02355v1. Full rights notice: licenses/arxiv-2605.02355-CC-BY-4.0.txt.\par\smallskip

\noindent\textit{Editorial scope.} Counted unit: the Lemma whose source label is \texttt{lem:crossing} in the article (source0.tex lines 1356--1358, source label \texttt{lem:crossing}), delivered together with the article's own complete original proof of it (source0.tex lines 1360--1426, one \texttt{proof} environment of 67 lines). No mathematics is written, repaired or re-derived: statement and proof are the article's own text line for line. Required context retained uncounted: lem:t-matching (lines 1321--1323); lem:t-matching-overlap (lines 1332--1337). Every definition, display and label of those lines is retained unchanged. Omitted from this package and not redistributed: 1--1320, 1324--1331, 1338--1355, 1359--1359, 1427--3134. Nothing inside a retained excerpt is omitted or altered: every retained body is delivered exactly as the article's lines are written, so no omission receipt of any kind accompanies this package. Citations: the retained excerpts carry no citation command at all, so no bibliography entry of the article is transcribed and no reference is redistributed. Reference adaptation: none was needed and none was made; every reference inside the delivered unit resolves inside it, so no unresolved reference or citation is produced. Printed slips kept literal: no printed word, symbol, hypothesis or number of the counted statement or of its proof was altered. Layout and class: the article's own class is \texttt{article} with packages including inputenc, fullpage, titlesec, authblk, xcolor, graphicx, subcaption, tikz, pgfplots, amssymb, dsfont, amsmath, amsthm, mathtools; the wrapper is the lane's pinned \texttt{amsart} editorial wrapper at 10pt on A4, which restates the theorem-like environments the retained text needs and adds the article's own macro definitions that the retained text needs (\texttt{\textbackslash Z}, \texttt{\textbackslash arr}, \texttt{\textbackslash dep}, \texttt{\textbackslash dotcup}, \texttt{\textbackslash energyact}, \texttt{\textbackslash w}, \texttt{\textbackslash wait}), with extra wrapper packages (bm, bbm, dsfont, xspace, cancel, mathtools, cleveref). The delivered numbering is the unit's own. Assets: no retained span contains a figure environment, a graphics-inclusion command, a PDF-page inclusion, a table or a TikZ picture; no image file is copied into this package, the package declares no asset dependency and no image file is redistributed with it. Licence: version arXiv:2605.02355v1 of this article is distributed under the Creative Commons Attribution 4.0 International licence, as recorded in the arXiv licence metadata for this exact version and shown on the abstract page https://arxiv.org/abs/2605.02355v1. A full rights notice is delivered at licenses/arxiv-2605.02355-CC-BY-4.0.txt. Characters: every retained excerpt is pure ASCII, so the delivered engine, XeLaTeX with its default Latin Modern text font, reports no missing character.
\begin{lemma}\label{lem:t-matching}
	The matching $\energyact_H \coloneqq \{ ((i, \dep),(i-1, \arr)) \mid i \in \{2, \dots , n \} \} \cup \{((1, \dep),(n, \arr)) \}$ is of maximum weight w.r.t.\ $\w$ among the matchings $\energyact$ that result in exactly one cycle $C = \energyact\dotcup A_{\wait}$.
\end{lemma}

\begin{lemma}
\label{lem:t-matching-overlap}
	The matching $\energyact_H$ from \Cref{lem:t-matching} yields a cycle $C = \energyact_H \dotcup A_{\wait}$ with maximum overlap 
	  \[o(C) = \sum_{i =1}^{n} t_i - t_{n} + \max\{t_{1} - \delta_{C},\, 0 \},\]
	where $t_1$ is the smallest acceleration/braking time of all $i \in C$ and $t_n$ is the largest.      
\end{lemma}

\begin{lemma}\label{lem:crossing}
	Let $C_1 \cup C_2 = \energyact\dotcup A_{\wait}$ be two crossing cycles obtained from a matching of energy arcs on a one-station network with $n$ trains. Then, the maximum overlap obtained by $\energyact$ is smaller or equal to the maximum overlap obtained by a maximum-weight Hamiltonian path matching $\energyact_H$, i.e. $o(\energyact) \leq o(\energyact_H)$.
\end{lemma}

\begin{proof}
	For the overlaps we obtain by \cref{lem:t-matching-overlap}
	\begin{align*}
		o(C_1) &= \sum_{i \in C_1} t_i - t_{l_1} + \max\{t_{s_1} - \delta_{C_1},\, 0 \}, \\
		o(C_2) &= \sum_{i \in C_2} t_i - t_{l_2} + \max\{t_{s_2} - \delta_{C_2},\, 0 \}, \\
		o(C) &= \sum_{i=1}^n t_i - t_n + \max\{t_1 - \delta_{C},\, 0 \}.
	\end{align*}
	This means
	\begin{align*}
		o(C_1) + o(C_2) = \sum_{i=1}^n t_i - t_{l_1} - t_{l_2} + \max\{t_{s_1} - \delta_{C_1},\, 0 \} + \max\{t_{s_2} - \delta_{C_2},\, 0 \}.
	\end{align*}
	We want to show that $o(C_1) + o(C_2) \leq o(C)$.
	
	\begin{itemize}
		\item If $t_{s_1} - \delta_{C_1} \leq 0$ and $t_{s_2} - \delta_{C_2} \leq 0$, then 
		\begin{align*}
			o(C_1) + o(C_2) = \sum_{i=1}^n t_i - t_{l_1} - t_{l_2} \leq \sum_{i=1}^n t_i - t_n \leq o(C),
		\end{align*}
		as one of $t_{l_1}$ and $t_{l_2}$ equals $t_n$.
		\item If $t_{s_1} - \delta_{C_1} \geq 0$ and $t_{s_2} - \delta_{C_2} \leq 0$, then 
		\begin{align*}
			o(C_1) + o(C_2) = \sum_{i=1}^n t_i - t_{l_1} - t_{l_2} + t_{s_1} - \delta_{C_1} \leq \sum_{i=1}^n t_i - t_n \leq o(C),
		\end{align*}
		because $\max\{t_{l_1}, t_{l_2}\} = t_n$, $\min\{t_{l_1}, t_{l_2}\} \geq t_{s_1}$, and $\delta_{C_1} \geq 0$.
		\item If $t_{s_1} - \delta_{C_1} \leq 0$ and $t_{s_2} - \delta_{C_2} \geq 0$, then 
		\begin{align*}
			o(C_1) + o(C_2) = \sum_{i=1}^n t_i - t_{l_1} - t_{l_2} + t_{s_2} - \delta_{C_2} \leq \sum_{i=1}^n t_i - t_n \leq o(C),
		\end{align*}
		because $\max\{t_{l_1}, t_{l_2}\} = t_n$, $\min\{t_{l_1}, t_{l_2}\} \geq t_{s_2}$ as $C_1$ and $C_2$ are crossing cycles, and $\delta_{C_2} \geq 0$.
		\item Finally, if $t_{s_1} - \delta_{C_1} \geq 0$ and $t_{s_2} - \delta_{C_2} \geq 0$, then
		\begin{align*}
			o(C_1) + o(C_2) &= \sum_{i=1}^n t_i - t_{l_1} - t_{l_2} + t_{s_1} - \delta_{C_1} + t_{s_2} - \delta_{C_2}\\ & = \sum_{i=1}^n t_i - t_{n} + t_{s_1} - (t^{\min}_{l} - t_{s_2}) - \delta_{C_1} - \delta_{C_2}  \\
            &= \sum_{i=1}^n t_i - t_{n} + t_1 - k  - \delta_{C_1} - \delta_{C_2},
		\end{align*}
		as it holds $t_{s_1} = t_1$ and $\max\{t_{l_1}, t_{l_2}\} = t_n$. Further, we define $t^{\min}_{l}:= \min\{t_{l_1}, t_{l_2}\} \geq t_{s_2}$, which holds as $C_1$ and $C_2$ are crossing cycles, and set $k \coloneqq t^{\min}_{l}- t_{s_2}$.
        If we consider the sum of the lower bounds and minima of the departure/arrival times on the energy arcs on cycles, we obtain
        \begin{align*}
            L_C &= L_{C_1} + L_{C_2} + k.
        \end{align*}

        Further, we know that that the sum of all lower and all upper bounds is independent of the matching. Hence, it holds
        \begin{align*}
            L_C + U_C &= L_{C_1} + U_{C_1} + L_{C_2} + U_{C_2}.
        \end{align*}
This yields
\begin{align*}
            U_C &= U_{C_1} + U_{C_2} - k.
        \end{align*}

        Hence, we get the following estimation for $\delta_C$:
        \begin{align*}
         \delta_C = d([L_C, U_C], T\Z) &=  d([L_{C_1} + L_{C_2} + k, U_{C_1} + U_{C_2} - k], T\Z) \\ 
        &\leq 
        d([L_{C_1} + L_{C_2}, U_{C_1} + U_{C_2}], T\Z) + k \\ 
        &\le  d([L_{C_1},\, U_{C_1}], T\Z) + d([L_{C_2}, U_{C_2}], T\Z) + k
        \\
        &= \delta_{C_1} + \delta_{C_2} + k.
        \end{align*}
		Therefore, we obtain the following result for the overlap of $C$:
		\begin{align*}
        o(C) &= \sum_{i=1}^{n} t_i - t_n + t_1 - \min\{t_1, \delta_C\}\\
        &\geq \sum_{i=1}^{n} t_i - t_n + t_1 - \min\{t_1, \delta_{C_1} + \delta_{C_2} + k\}\\
        &\geq \sum_{i=1}^{n} t_i - t_n + t_1  - k - \delta_{C_1} - \delta_{C_2} \\
        &= o(C_1) + o(C_2).
		\end{align*}
	\end{itemize}
\end{proof}

\end{document}
